\section{Usage Guidelines}
\label{sec:guidelines}

In addition to the DPV-ODRL mapping
described in the previous section,
we provide a guidance document 
at \url{https://dev.dpvcg.org/guides/dpv-odrl}
with examples of ODRL policies that use the developed mapping
and conform with the ODRL standard.
In this section,
we reproduce a few of these examples to 
showcase the usefulness of this work
to tackle legal compliance enforcement
in data spaces.

Listing~\ref{lst:partycollection}
provides an example of a DPV-ODRL policy,
where a data holder allows the usage of a certain asset,
i.e., \texttt{ex:research-dataset},
only for the purpose of
\texttt{dpv:ResearchAndDevelopment}
and only for data users which are classified as 
\texttt{dpv:AcademicScientificOrganisation}.

\begin{scriptsize}
\begin{lstlisting}[language=turtle, caption=Purpose-restricted agreement with a party collection refinement., label=lst:partycollection]
ex:policy-partycollection a odrl:Agreement ;
    odrl:uid ex:policy-partycollection ;
    odrl:profile dpv-odrl: ;
    odrl:permission [
        odrl:assigner ex:data-holder ;
        odrl:assignee [ 
            a odrl:PartyCollection ;
            odrl:source ex:entities ;
            odrl:refinement [
                odrl:leftOperand dpv-odrl:Entity ;
                odrl:operator odrl:isA ;
                odrl:rightOperand dpv:AcademicScientificOrganisation ] ] ;
        odrl:target ex:research-dataset ;
        odrl:action dpv-odrl:Use ;
        odrl:constraint [
            odrl:leftOperand dpv-odrl:Purpose ;
            odrl:operator odrl:isA ;
            odrl:rightOperand dpv:ResearchAndDevelopment ] ] .
\end{lstlisting}
\end{scriptsize}

Listing~\ref{lst:assetcollection}
provides an example of a DPV-ODRL policy,
where a data subject allows the usage of elements
of a certain asset collection,
i.e., \texttt{ex:data-catalog},
that are instances of health record data,
i.e., \texttt{pd:HealthRecord}.

\begin{scriptsize}
\begin{lstlisting}[language=turtle, caption=Policy with a DPV-ODRL-based asset collection refinement on personal data type., label=lst:assetcollection]
ex:policy-assetcollection a odrl:Policy ;
    odrl:uid ex:policy-assetcollection ;
    odrl:profile dpv-odrl: ;
    odrl:permission [
        odrl:target [ 
            a odrl:AssetCollection ;
            odrl:source ex:data-catalog ;
            odrl:refinement [
                odrl:leftOperand dpv-odrl:PersonalData ;
                odrl:operator odrl:isA ;
                odrl:rightOperand pd:HealthRecord ]] ;
        odrl:assigner ex:data-subject ;
        odrl:action dpv-odrl:Use ] .
\end{lstlisting}
\end{scriptsize}

Listing~\ref{lst:righterasure} contains an
example specifying an agreement between
a data subject and
a data controller,
where the data controller has the obligation
to delete the data subject data,
given that the latter exercised their
GDPR right to erasure, i.e., \texttt{eu-gdpr:A17},
supported by a non-necessity justification,
i.e., \texttt{eu-gdpr:JustificationA17NonNecessity},
which can be used in such a request
when the involved personal data is no longer necessary
in relation to the purposes for which it was initially collected
or otherwise processed.

\begin{scriptsize}
\begin{lstlisting}[language=turtle, caption=Obligation-based policy for the exercising of a right and related justification., label=lst:righterasure]
ex:policy-right-erasure-justification a odrl:Agreement ;
    odrl:uid ex:policy-right-erasure-justification ;
    odrl:profile dpv-odrl: ;
    odrl:obligation [
        odrl:assigner ex:subject ;
        odrl:assignee ex:controller ;
        odrl:target ex:subject-data ;
        odrl:action dpv-odrl:Delete ;
        odrl:constraint [
            odrl:leftOperand dpv-odrl:Right ;
            odrl:operator odrl:eq ;
            odrl:rightOperand eu-gdpr:A17 ], [
            odrl:leftOperand dpv-odrl:Justification ;
            odrl:operator odrl:eq ;
            odrl:rightOperand eu-gdpr:JustificationA17NonNecessity ] ] .
\end{lstlisting}
\end{scriptsize}

Finally,
we provide a policy catalog at
\url{https://github.com/besteves4/SDS2026-DPV-ODRL-mapping}
that contains an example policy for each mapped concept.

\subsection{SHACL Shapes for Constraint Validation}
\label{sec:shacl}

Beyond specifying examples,
the guidance document
recommends the usage of SHACL\footnote{
    \url{https://www.w3.org/TR/shacl/}
} shapes
to validate DPV-ODRL-based policies.
While such shapes already exist
to validate core ODRL constructs\footnote{See \url{
    https://github.com/oeg-upm/licensius/blob/master/odrlapi/src/main/resources/shapes.ttl}},
the same exercise is needed for
the policies described in this work,
in particular to validate the usage of constraints,
as previously discussed in Section~\ref{sec:operators}.
In this context,
the ODRL vocabulary has a set of 12
defined constraint operators to express
the relation between left and right operands.
Given that multiple operators
can be used to restrict this relation,
the left operand and operator usage
should be refined for each left operand,
e.g., the \texttt{odrl:isA} operator
should not be used with the
\texttt{dpv-odrl:Frequency} left operand.
Listing~\ref{lst:shacl}
provides a template 
for a SPARQL-based SHACL constraint
that can be used to
detect the usage of an invalid operator
for a given left operand.

\pagebreak

\begin{scriptsize}
\begin{lstlisting}[language=turtle, caption=Template shape to detect the usage of an invalid operator for a given left operand., label=lst:shacl]
ex:ConstraintShapeTemplate a sh:NodeShape ;
    sh:targetClass odrl:Constraint ;
    sh:sparql [
        a sh:SPARQLConstraint ;
        sh:message "Invalid operator for given left operand." ;
        sh:prefixes odrl: ;
        sh:select """
            SELECT $this WHERE {
                $this odrl:leftOperand <LEFT_OPERAND> .
                $this odrl:operator ?op .
                FILTER(?op NOT IN (<ALLOWED_OPERATORS>))
            }
        """ ;
    ] .
\end{lstlisting}
\end{scriptsize}